% Nouveau EGA V (cote 25): archive cote 25, site page 46 (PDF page 47).
% typescript. Draft compared: none.
% AI-transcribed by Claude Opus 5.5 (claude-opus-5-5, adaptive thinking, effort high), Message Batches API, batch ega25b (msgbatch_01A7n5YuXubUPNK1zsepi8RG), 2026-11-13; not yet checked by eye.
\scanpage{25/46}
\mspage{42}

teur de X en un point $x\in X(k)$, \xout{\unsure{mais}} tel qu'il existe un \uline{autre point}
$x'\in X(k)$ en lequel le même hyperplan soit tangent en x. On en fait des
exemples évidents avec \xout{\unsure{des exemples}} P de dimension 2, X une courbe non si
singulière de degré $\geqslant$ \unsure{4} , en c aractéristique quelconque. \inserted{[NB les ``tangen-tes doubles'' = \unsure{X} correspondant aux pts doubles de la ``courbe duale'' !].}

\uline{Corollaire} 8.\over{\unsure{10}}{11}. Supposons k de caractéristique nulle. Alors:

a) L'image de $\underline{Y}^{\text{sup sing}}$ dans $\overset{\hand{\vee}}{P}$ est de codimension $\geqslant 2$.

b) La condition (iii) de 8.\over{\unsure{7}}{8}. équivaut aux autres conditions, i.e.
la négation des autres conditions\hand{,} \inserted{sauf $\underline{Y}^{\text{sing}}=\underline{Y}^{\text{sup sing}}$,} signifie aussi que l'image de $\underline{Y}^{\text{sing}}$
dans $\overset{\hand{\vee}}{P}$ est de codimension $\geqslant 2$.

Evidemment l'assertion b) implique a), compte tenu de 8.\over{\unsure{7}}{8}. Or
par la théorie de la dimension, (iii) signifie que L/K est une extensi-%
on finie, \xout{\unsure{et en}} (on pourrait le mettre sous forme (iii bis) dans 8.\over{\unsure{7}}{8}.),
et en car 0 \xout{\unsure{cela implique}} L est de toutes façons séparable sur K, d'où
la condition (i bis) de 8.\over{\unsure{7}}{8}.

\uline{Remarques} 8.1\over{1}{2}.. a) Géométriquement, l'assertion a) signifie essen-%
tiellement que pour un pinceau linéaire assez général de sections hyper-%
planes, tous les membres du pinceau sont lisses ou \xout{n}ont \xout{\unsure{que des points}}
\xout{\unsure{gé}} pour s euls points géométriquement des points géométriquement
singuliers ordinaires (et en f ait, comme on constate aussitôt - et il
y aurait lieu d'énoncer a) en conséquence sous une forme un peu plus
précise - ont au plus \uline{un} tel point géométriquement singulier ). L'as\xout{s}-%
sertion b) signifie \inserted{essentiellement} que lorsque $\underline{Y}^{\text{sing}}=\underline{Y}^{\text{sup sing}}$ (ce qui s'exprime
analytiquement par l\xout{\unsure{'annulation}} \inserted[typed]{nullité} d'une certaine section D d'un Module
inversible $\underline{\hand{\omega}}^{\otimes 2}_{X/\hand{k}}\otimes\underline{O}_{\underline{Y}^1}(1,1)$ sur $\underline{Y}^1$), alors pour un pinceau linéaire
assez général de sections hyperplanes, tous les membres de ce \xout{\unsure{faisceau}}
pinceau sont lisses. Cette deuxième situation (qu'on s\xout{x}oit d'ailleurs
en car.0 ou non ) doit être considérée comme tout à fait exceptionnelle

\mspage{46}
